%
% architecture.tex
%
% Dungeon: Technical Architecture Specification
% Written for engineers implementing the system.
%
% Compile: pdflatex architecture.tex
%

\documentclass[a4paper,10pt]{article}
\usepackage[margin=1in]{geometry}
\usepackage{booktabs}
\usepackage{listings}
\usepackage{xcolor}
\usepackage{enumitem}
\usepackage{longtable}
\usepackage[colorlinks=true,linkcolor=blue,citecolor=blue,urlcolor=blue]{hyperref}

\lstset{
  basicstyle=\ttfamily\small,
  keywordstyle=\color{blue!80!black},
  commentstyle=\color{green!50!black},
  stringstyle=\color{red!60!black},
  breaklines=true,
  frame=single,
  framerule=0.3pt,
  rulecolor=\color{gray!50},
  backgroundcolor=\color{gray!5},
  columns=fullflexible,
}

\newcommand{\iface}[1]{\texttt{#1}}
\newcommand{\tool}[1]{\texttt{#1}}
\newcommand{\cmd}[1]{\texttt{#1}}
\newcommand{\field}[1]{\texttt{#1}}
\newcommand{\mode}[1]{\textbf{#1}}
\newcommand{\pkg}[1]{\texttt{#1}}

\begin{document}

\title{Dungeon: Technical Architecture Specification}
\author{LLM-Directed Text Adventure with Go Game Engine}
\date{March 2026 --- v0.2 (design-stage)}
\hypersetup{
  pdftitle={Dungeon: Technical Architecture Specification},
  pdfauthor={Punt Labs},
  pdfsubject={Architecture},
  pdfcreator={pdflatex}
}
\maketitle

\tableofcontents
\newpage

% ===================================================================
\section{Overview}
% ===================================================================

Dungeon is a text adventure game for Claude Code.  In its current
form the entire engine --- parser, rules, narrator, state machine ---
is a single LLM skill prompt.  This specification describes the
evolved architecture: a deterministic Go game engine that handles
all state transitions, with the LLM acting as Dungeon Master
(narrator, cartographer, encounter designer) rather than parser.

The central insight is \textbf{play mode as a first-class concept}.
A play mode is a named combination of three swappable Go interfaces:
\iface{CommandInterpreter}, \iface{Narrator}, and \iface{Renderer}.
All modes share the same engine, the same save files, and the same
scenario format.  Switching modes mid-adventure --- save in
\mode{dm}, resume in \mode{solo} --- is explicitly supported.

\subsection{Mechanics Reference}

The game targets the mechanical vocabulary of \emph{Wizardry~I}
(grid-based dungeon, turn-based combat, four character classes,
dungeon floors with stairs) and the exploratory sensibility of
\emph{Zork~I} (free-text examination, rich environmental prose,
consequence of curiosity).  Where the two conflict, Wizardry's
determinism wins: rules are authoritative, the engine adjudicates.

\subsection{Design Principles}

\begin{enumerate}[label=\textbf{P\arabic*}.]
  \item \textbf{The engine adjudicates; the LLM narrates.}
    No game outcome is left to probabilistic language model output.
    State transitions, combat resolution, and inventory rules are
    computed deterministically in Go.  The LLM enriches the result
    with prose.
  \item \textbf{Interfaces are the seams.}
    \iface{CommandInterpreter}, \iface{Narrator}, and
    \iface{Renderer} are the only coupling points between the engine
    and external systems.  Each can be replaced independently.
  \item \textbf{Play mode is operational context, not architecture.}
    The same binary supports DM, solo, and headless play.  No code
    paths are conditional on mode; mode selects implementations.
  \item \textbf{Daemon scope stays minimal.}
    The shared game engine daemon handles game logic and push
    routing only.  No LLM calls, no orchestration, no business
    logic outside the game model.
  \item \textbf{Save files are cross-mode.}
    A save created in \mode{dm} mode is loadable in \mode{solo} or
    \mode{headless}.  The \field{play\_mode} field in the save is
    advisory; \cmd{--mode} overrides it at resume time.
  \item \textbf{Party is designed in, not bolted on.}
    \texttt{GameState.Party} is always a slice (length~1 for single
    player).  Combat, movement, and save tooling accept an optional
    \field{character\_id}.  Multi-player (Biff Phase~5) requires no
    engine redesign.
\end{enumerate}

% ===================================================================
\section{Play Modes}
% ===================================================================

Three play modes are first-class citizens of the system.  Each is a
named triple of interface implementations.

\begin{center}
\begin{tabular}{llll}
\toprule
\textbf{Interface} & \textbf{dm} & \textbf{solo} & \textbf{headless} \\
\midrule
\iface{CommandInterpreter}
  & \texttt{LLMInterpreter}
  & \texttt{SLMInterpreter}
  & \texttt{SLMInterpreter} or \texttt{RulesInterpreter} \\
\iface{Narrator}
  & \texttt{LLMNarrator}
  & \texttt{SLMNarrator}
  & \texttt{TemplateNarrator} \\
\iface{Renderer}
  & \texttt{LuxRenderer}
  & \texttt{LuxRenderer} or \texttt{CLIRenderer}
  & \texttt{CLIRenderer} \\
\textbf{Engine deployment}
  & daemon
  & embedded
  & embedded \\
\midrule
\textbf{Invoked by}
  & \cmd{/dungeon} skill
  & \cmd{dungeon solo}
  & \cmd{dungeon headless} \\
\bottomrule
\end{tabular}
\end{center}

\subsection{Mode Responsibilities}

\begin{center}
\begin{tabular}{lp{2.8cm}p{2.8cm}p{2.8cm}}
\toprule
\textbf{Concern} & \textbf{dm} & \textbf{solo} & \textbf{headless} \\
\midrule
Room narration        & LLM             & SLM             & Template \\
Free-text parsing     & LLM             & SLM             & Rules/regex \\
Scenario generation   & LLM interactive & Pre-made only   & Pre-made only \\
Encounter lore        & LLM             & Pre-made only   & Pre-made only \\
State transitions     & Engine          & Engine          & Engine \\
Combat resolution     & Engine          & Engine          & Engine \\
Fog-of-war, save/load & Engine          & Engine          & Engine \\
Rendering             & Lux             & Lux or CLI      & CLI \\
\bottomrule
\end{tabular}
\end{center}

% ===================================================================
\section{System Architecture}
% ===================================================================

\subsection{Component Map}

\begin{center}
\begin{tabular}{lp{3cm}p{7cm}}
\toprule
\textbf{Component} & \textbf{Language / Layer} & \textbf{Responsibility} \\
\midrule
Game engine
  & Go (L1)
  & All state transitions; character, map, inventory, combat, save/load.
    Zero LLM dependencies. \\
\iface{CommandInterpreter}
  & Go interface
  & Maps player free-text to a structured \texttt{EngineAction}. \\
\iface{Narrator}
  & Go interface
  & Generates human-readable prose from engine event results. \\
\iface{Renderer}
  & Go interface
  & Presents game state to the player (Lux HUD or ANSI CLI). \\
MCP server
  & Go binary (\cmd{dungeon serve})
  & Exposes game engine as MCP tools; NDJSON over Unix socket in
    daemon mode; stdio in embedded mode. \\
SKILL.md
  & Prompt (L4)
  & DM role instructions for Claude; narration, scenario generation,
    lore, examination responses. \\
mcp-proxy
  & Go shim (design-stage)
  & Bridges Claude Code stdio MCP to shared daemon; injects session
    identity; enables multi-player push. \\
Lux
  & Python / ImGui
  & Native display surface; receives element trees from
    \texttt{LuxRenderer} via MCP tools. \\
ollama
  & External service
  & Serves SLM for \texttt{SLMInterpreter} and \texttt{SLMNarrator}
    in \mode{solo}/\mode{headless} modes. \\
\bottomrule
\end{tabular}
\end{center}

\subsection{Topology: Embedded vs.\ Daemon}

The engine runs in one of two deployment topologies depending on
play mode.  The MCP tool API and all interface implementations are
identical in both topologies.

\begin{lstlisting}[title={Embedded (solo / headless)}]
player (stdin / Lux)
    |
dungeon solo            <-- single process
    |-- Go engine       embedded, no socket
    |-- CommandInterpreter (SLM or rules)
    |-- Narrator        (SLM or template)
    `-- Renderer        (Lux MCP calls / ANSI stdout)
\end{lstlisting}

\begin{lstlisting}[title={Daemon (dm mode / multi-player)}]
Claude Code session (DM)
    `-- mcp-proxy  <10ms, <10MB
            |
            | NDJSON Unix socket
            v
    dungeon serve   <-- shared daemon, one per adventure
            |-- GameState (all sessions)
            |-- session identity routing
            `-- tools/list_changed push to DM on player action

Claude Code session (player, future)
    `-- mcp-proxy -----^
\end{lstlisting}

The daemon holds all game state.  It knows which connection is the
DM versus a player via session identity injected by mcp-proxy
(resolved from the topmost \cmd{claude} ancestor PID at connect
time).  DM-privileged tools are only available to the DM session.

\subsection{Data Flow: DM Mode}

\begin{enumerate}
  \item \textbf{Button press (navigation/combat).}
    Player clicks a Lux navigation button.  Lux emits an interaction
    event.  \texttt{LuxRenderer} calls the engine directly:
    \tool{move("N")} or \tool{attack(target\_id)}.
    The engine updates state and returns the result.
    \texttt{LuxRenderer} patches the HUD (map fog reveal, HP update)
    via \tool{lux.update()}.
    \texttt{LLMNarrator} is called with the result to generate
    flavour text, which is appended to the narration log.
    \emph{No LLM round-trip occurs for the state transition itself.}
  \item \textbf{Free-text input.}
    Player types ``I search the walls for a hidden door.''
    \texttt{LLMInterpreter} sends the text to Claude with the current
    room context.  Claude returns a structured action:
    \texttt{\{action: "search", target: "walls", hint: "door"\}}.
    Engine executes \tool{search(target, hint)}.
    Result (found/not found, item if any) passed to
    \texttt{LLMNarrator}.  Claude generates flavour prose.
    \texttt{LuxRenderer} updates the narration log.
  \item \textbf{State change push (multi-player, future).}
    After any player action, the daemon emits
    \texttt{tools/list\_changed} to the DM session.
    Claude Code (DM) notices the notification and calls
    \tool{get\_game\_state} to read the delta.
    DM narrates the result to all player sessions.
\end{enumerate}

\subsection{Data Flow: Solo Mode}

Solo mode is functionally identical to DM mode with two substitutions:
\texttt{LLMInterpreter} $\to$ \texttt{SLMInterpreter} (ollama HTTP,
port 11434) and \texttt{LLMNarrator} $\to$ \texttt{SLMNarrator}.
The engine runs embedded; there is no daemon or proxy.  If ollama is
unavailable, \texttt{SLMInterpreter} degrades to
\texttt{RulesInterpreter} automatically.

% ===================================================================
\section{The Three Engine Interfaces}
\label{sec:interfaces}
% ===================================================================

\subsection{CommandInterpreter}

Maps a player's free-text string to a structured
\texttt{EngineAction} that the engine can execute deterministically.

\begin{lstlisting}[language=Go]
type CommandInterpreter interface {
    Interpret(ctx context.Context, input string,
              state GameState) (EngineAction, error)
}
\end{lstlisting}

\begin{center}
\begin{tabular}{lll}
\toprule
\textbf{Implementation} & \textbf{Used in} & \textbf{Mechanism} \\
\midrule
\texttt{LLMInterpreter}
  & \mode{dm}
  & MCP callback to Claude; compact structured prompt; returns JSON action. \\
\texttt{SLMInterpreter}
  & \mode{solo}, \mode{headless}
  & ollama HTTP API; phi-3-mini (default); verb/noun classification. \\
\texttt{RulesInterpreter}
  & \mode{headless} fallback
  & Regex/keyword matching over known verb--noun pairs; no model required. \\
\bottomrule
\end{tabular}
\end{center}

\subsection{Narrator}

Generates human-readable prose from an \texttt{EngineEvent}: the
result of an action after the engine has resolved it.

\begin{lstlisting}[language=Go]
type Narrator interface {
    Narrate(ctx context.Context, event EngineEvent,
            state GameState) (string, error)
}
\end{lstlisting}

\begin{center}
\begin{tabular}{lll}
\toprule
\textbf{Implementation} & \textbf{Used in} & \textbf{Mechanism} \\
\midrule
\texttt{LLMNarrator}
  & \mode{dm}
  & Claude generates rich contextual prose; full adventure context provided. \\
\texttt{SLMNarrator}
  & \mode{solo}
  & ollama generates short atmospheric text from room seed + event type. \\
\texttt{TemplateNarrator}
  & \mode{headless}
  & Go text/template: ``You move north. You are in the \{room.name\}.'' \\
\bottomrule
\end{tabular}
\end{center}

\subsection{Renderer}

Presents the current \texttt{GameState} to the player after each
turn.  Also receives button interaction events in DM/solo modes.

\begin{lstlisting}[language=Go]
type Renderer interface {
    Render(ctx context.Context, state GameState,
           narration string) error
    Events() <-chan InputEvent  // button presses, free-text
}
\end{lstlisting}

\begin{center}
\begin{tabular}{lll}
\toprule
\textbf{Implementation} & \textbf{Used in} & \textbf{Mechanism} \\
\midrule
\texttt{LuxRenderer}
  & \mode{dm}, \mode{solo}
  & Calls Lux MCP tools (\tool{show}, \tool{update}, \tool{recv})
    to drive the full HUD. \\
\texttt{CLIRenderer}
  & \mode{headless}, SSH
  & ANSI terminal; ASCII map, text HP bar, stdin input loop. \\
\bottomrule
\end{tabular}
\end{center}

% ===================================================================
\section{Game Engine Design}
\label{sec:engine}
% ===================================================================

\subsection{Core Data Model}

\begin{lstlisting}[language=Go]
type GameState struct {
    ScenarioID  string
    PlayMode    string   // advisory; --mode overrides on load
    Timestamp   string
    Character   Character
    Party       []Character  // len 1 for single-player
    Dungeon     DungeonState
    Log         []LogEntry
}

type Character struct {
    Name       string
    Class      CharClass   // Fighter | Mage | Thief | Priest
    Level      int
    HP, MaxHP  int
    MP, MaxMP  int
    Stats      Stats       // STR INT DEX CON WIS CHA
    XP, Gold   int
    Inventory  []Item
    Equipped   Equipment
    Conditions []Condition
}

type DungeonState struct {
    CurrentRoom  string
    Visited      []string
    RoomState    map[string]RoomState
}
\end{lstlisting}

\subsection{Character}

\begin{itemize}
  \item \textbf{Classes}: Fighter, Mage, Thief, Priest.
    Class determines HP dice, MP pool, available spells, and XP
    thresholds.
  \item \textbf{Stats}: STR, INT, DEX, CON, WIS, CHA.
    Bonuses computed by standard table (8--18 range, D\&D-adjacent).
  \item \textbf{Equipment slots}: weapon, armour, ring, amulet.
  \item \textbf{Leveling}: XP thresholds defined per class in
    scenario YAML.  Stat gains on level-up are deterministic.
\end{itemize}

\subsection{Map}

\begin{itemize}
  \item Grid-based dungeon in the Wizardry~I style (20$\times$20
    floor).  Rooms are nodes; corridors are edges.
  \item \textbf{Pre-generated}: the DM creates the room graph before
    play begins, either by writing scenario YAML directly or via the
    \cmd{/dungeon:create} skill command.  Procedural generation is
    explicitly not used --- it produces maps that are inconsistent on
    backtracking.
  \item \textbf{Fog of war}: unvisited rooms are opaque on the Lux
    map canvas.  Room content is not revealed until the player enters.
  \item \textbf{Room types}: corridor, chamber, stairway (up/down),
    treasure vault, boss chamber.
  \item \textbf{Door kinds}: open, locked (requires key item),
    hidden (requires search action), one-way, trapped.
\end{itemize}

\subsection{Combat}

Turn-based, deterministic.  No probabilities are left to the LLM.

\begin{itemize}
  \item Initiative roll (DEX modifier + 1d10) determines turn order.
  \item \textbf{Player actions per turn}: Attack, Cast Spell, Use
    Item, Flee (DEX check), Defend (+AC until next turn).
  \item \textbf{Enemy AI}: patterns encoded in scenario YAML.
    Supported patterns: \texttt{aggressive} (always attack),
    \texttt{cautious} (flee below 30\% HP), \texttt{support}
    (buff allies before attacking), \texttt{scripted} (fixed action
    sequence).
  \item Damage formula: weapon dice + STR modifier vs.\ target AC.
  \item Conditions: poisoned, stunned, confused, blessed, hasted.
    Duration in turns; applied/cleared by spells and items.
  \item XP is awarded at \texttt{end\_combat}; loot rolls are
    seeded from scenario YAML drop tables.
\end{itemize}

% ===================================================================
\section{Data Formats}
% ===================================================================

\subsection{Scenario Format (YAML)}

Scenarios are pre-authored YAML files stored in \texttt{scenarios/}.
The engine parses them at startup; the LLM never sees the raw YAML
during play.

\begin{lstlisting}[title={scenarios/classic-fantasy-dungeon.yaml (excerpt)}]
id: classic-fantasy-dungeon
title: "The Grimhold Dungeon"
description: "A classic dungeon crawl..."
version: "1.0"
xp_thresholds:
  fighter: [0, 2000, 4000, 8000, 16000]
  mage:    [0, 2500, 5000, 10000, 20000]

rooms:
  - id: entrance
    name: "Dungeon Entrance"
    description: "Stone stairs descend into darkness."
    exits:
      south: goblin_lair
    items: []
    enemies: []

  - id: goblin_lair
    name: "Goblin Lair"
    description: "The reek of smoke and old blood."
    exits:
      north: entrance
      east:  { room: treasure_vault, door: locked, key: rusty_key }
    enemies:
      - template: goblin
        count: 3
    loot_table:
      - item: rusty_key
        drop_chance: 0.8

enemy_templates:
  goblin:
    hp_dice: "2d6"
    ac: 12
    attack_dice: "1d4"
    xp: 30
    ai: aggressive
\end{lstlisting}

\subsection{Save File Format (JSON)}

Save files live at \texttt{.dungeon/saves/<slot>.json} (gitignored).

\begin{lstlisting}[title={.dungeon/saves/1.json (excerpt)}]
{
  "version": 1,
  "scenario": "classic-fantasy-dungeon",
  "play_mode": "dm",
  "timestamp": "2026-03-10T18:00:00Z",
  "character": {
    "name": "Aldric", "class": "fighter", "level": 3,
    "hp": 45, "max_hp": 60, "mp": 0, "max_mp": 0,
    "stats": {"str":16,"dex":12,"con":14,"int":9,"wis":10,"cha":11},
    "xp": 1240, "gold": 48,
    "inventory": [],
    "equipped": {"weapon": "rusty_sword"},
    "conditions": []
  },
  "party": [],
  "dungeon": {
    "current_room": "goblin_lair",
    "visited_rooms": ["entrance", "corridor_a", "goblin_lair"],
    "room_state": {
      "goblin_lair":    {"cleared": true, "items_taken": ["rusty_key"]},
      "treasure_vault": {"door_locked": false}
    }
  },
  "adventure_log": []
}
\end{lstlisting}

\texttt{party} is always present (empty for single-player).
\texttt{play\_mode} is advisory; \cmd{--mode} overrides on resume.

% ===================================================================
\section{MCP Tool Surface (Game Engine)}
% ===================================================================

The engine exposes all game actions as MCP tools.  The interface is
identical in embedded and daemon topologies.

\subsection{Session Management}

\begin{center}
\begin{tabular}{lp{4cm}p{5cm}}
\toprule
\textbf{Tool} & \textbf{Arguments} & \textbf{Returns} \\
\midrule
\tool{new\_game}
  & \field{scenario\_id}, \field{char\_name}, \field{char\_class},
    \field{mode?}
  & Full initial game state \\
\tool{load\_game}
  & \field{slot?}, \field{mode?}
  & Full game state from save \\
\tool{save\_game}
  & \field{slot?}
  & Save file path \\
\tool{list\_saves}
  & ---
  & \texttt{[\{slot, scenario, char, mode, timestamp\}]} \\
\tool{list\_scenarios}
  & ---
  & \texttt{[\{id, title, description\}]} \\
\bottomrule
\end{tabular}
\end{center}

\subsection{Navigation}

\begin{center}
\begin{tabular}{lp{4cm}p{5cm}}
\toprule
\textbf{Tool} & \textbf{Arguments} & \textbf{Returns} \\
\midrule
\tool{move}
  & \field{direction}, \field{character\_id?}
  & new\_room, encounter?, items?, event? \\
\tool{look}
  & ---
  & Current room exits, items, enemies \\
\tool{get\_map}
  & ---
  & Room graph with visited/fog state \\
\bottomrule
\end{tabular}
\end{center}

\subsection{Character and Inventory}

\begin{center}
\begin{tabular}{lp{4cm}p{5cm}}
\toprule
\textbf{Tool} & \textbf{Arguments} & \textbf{Returns} \\
\midrule
\tool{get\_character} & --- & Stats, HP/MP, level, XP, gold, conditions \\
\tool{get\_inventory} & --- & Item list with weights and equipped slots \\
\tool{pick\_up}   & \field{item\_id}  & success, weight\_remaining \\
\tool{drop}       & \field{item\_id}  & success \\
\tool{equip}      & \field{item\_id}  & success, stat\_delta \\
\tool{use\_item}  & \field{item\_id}, \field{target?} & Effect result \\
\tool{examine}    & \field{target}    & Raw detail (Narrator adds lore) \\
\bottomrule
\end{tabular}
\end{center}

\subsection{Combat}

\begin{center}
\begin{tabular}{lp{4cm}p{5cm}}
\toprule
\textbf{Tool} & \textbf{Arguments} & \textbf{Returns} \\
\midrule
\tool{get\_combat\_state}
  & ---
  & Turn order, HP, conditions \\
\tool{attack}
  & \field{target\_id}, \field{weapon?}, \field{character\_id?}
  & Roll, damage, target HP, outcome \\
\tool{cast\_spell}
  & \field{spell\_id}, \field{target?}, \field{character\_id?}
  & Effect, MP cost, outcome \\
\tool{defend}
  & \field{character\_id?}
  & Defence bonus until next turn \\
\tool{use\_item\_combat}
  & \field{item\_id}, \field{target?}, \field{character\_id?}
  & Effect result \\
\tool{flee}
  & \field{character\_id?}
  & Success (DEX check); penalty on fail \\
\tool{end\_combat}
  & ---
  & XP gained, loot found \\
\bottomrule
\end{tabular}
\end{center}

\subsection{Search and Interaction}

\begin{center}
\begin{tabular}{lp{4cm}p{5cm}}
\toprule
\textbf{Tool} & \textbf{Arguments} & \textbf{Returns} \\
\midrule
\tool{search}
  & \field{target?}, \field{hint?}
  & Found items/secrets, or nothing \\
\tool{interact}
  & \field{object\_id}, \field{action?}
  & Outcome (lever pulled, door opened\ldots) \\
\tool{rest}
  & ---
  & HP/MP restored; random encounter check \\
\bottomrule
\end{tabular}
\end{center}

% ===================================================================
\section{Lux Display Integration}
% ===================================================================

\subsection{HUD Layout}

The \texttt{LuxRenderer} drives a four-panel HUD.  Navigation and
combat buttons deliver input events directly to the engine --- no
LLM round-trip is required for button-driven actions.

\begin{lstlisting}[title={Lux HUD layout (schematic)}]
+--------------------------------------------------+
| Menu: [Game] [Character] [Map]                   |
+----------------------------+---------------------+
|                            | HP  [========--] 45 |
|   DUNGEON MAP              | MP  [----------]  0 |
|   (draw canvas,            | Fighter Lv3  XP 1k  |
|    grid rooms,             +---------------------+
|    fog of war,             | EXITS               |
|    visited/unvisited,      |       [N]            |
|    player dot,             |  [W]       [E]      |
|    locked doors)           |       [S]            |
|                            |       [D stairs]     |
+----------------------------+---------------------+
| NARRATION LOG              | COMBAT (when active) |
| (markdown, scrolling)      | Goblin [=====--] 8hp |
|                            | [Attack] [Use Item]  |
|                            | [Defend] [Flee]      |
|                            | Round 2 - Your turn  |
+----------------------------+---------------------+
\end{lstlisting}

\subsection{Update Strategy}

\begin{center}
\begin{tabular}{lp{8cm}}
\toprule
\textbf{Event} & \textbf{Lux operation} \\
\midrule
Room transition (new room entered)
  & \tool{lux.show()} --- replace full narration panel, update map fog \\
HP/MP change during combat
  & \tool{lux.update()} --- patch progress bar elements in-place \\
Narration text appended
  & \tool{lux.update()} --- append to narration log list \\
Combat begins/ends
  & \tool{lux.show()} --- activate or deactivate combat panel \\
Map fog reveal (new room)
  & \tool{lux.update()} --- patch map canvas element \\
Nav button press (player)
  & \tool{lux.recv()} event $\to$ engine directly \\
\bottomrule
\end{tabular}
\end{center}

\subsection{Lux Tools Used}

The \texttt{LuxRenderer} uses the following Lux MCP tools:

\begin{itemize}
  \item \tool{show(scene\_id, elements, title?)} --- initial HUD render
    and full scene transitions.
  \item \tool{update(scene\_id, patches)} --- incremental HP, MP, fog,
    and log updates within a scene.
  \item \tool{recv(timeout?)} --- poll for button interaction events
    (navigation presses, combat actions).
\end{itemize}

The engine does not use \tool{show\_table}, \tool{show\_dashboard},
or \tool{show\_diagram}; those are data-analysis surfaces.

% ===================================================================
\section{mcp-proxy Integration}
% ===================================================================

\textbf{Status}: design-stage.  The mcp-proxy binary does not yet
exist.  DM mode (Phase~2) should initially use direct stdio MCP and
migrate to the proxy pattern when mcp-proxy ships.

\subsection{Problem}

Claude Code spawns one MCP server process per session via stdio.
In DM mode, that process is the game engine daemon.  If the DM's
Claude session owns the engine process via stdio, a second player
connection is architecturally impossible: there is no shared process
to connect to, and the DM cannot push events to players.

\subsection{Solution}

mcp-proxy is a near-zero-cost Go shim ($<$10~ms startup,
$<$10~MB RSS) that each Claude session spawns via stdio.  The shim
connects to a pre-running shared daemon over a persistent
bidirectional transport (WebSocket or NDJSON Unix socket) and
relays MCP messages.  At connect time, mcp-proxy resolves the
topmost \cmd{claude} ancestor PID and injects it as a session
identity header.

The daemon uses session identity to:

\begin{itemize}
  \item Route DM-privileged tools only to the DM session.
  \item Push \texttt{tools/list\_changed} to the DM when a player
    takes an action.
  \item Route DM narration output to all player sessions.
\end{itemize}

\subsection{Daemon Scope Constraint}

The beads project accumulated 70\,000 lines of daemon code before
being deleted.  This daemon must not repeat that mistake.  The
game daemon handles:

\begin{itemize}
  \item Game state (authoritative \texttt{GameState} for the session).
  \item MCP tool dispatch (receive call $\to$ engine method $\to$
    return result).
  \item Push routing (\texttt{tools/list\_changed}, narration
    broadcast).
\end{itemize}

It does not handle: LLM calls, session management, orchestration
logic, plugin lifecycle, or anything outside the game model.

% ===================================================================
\section{SLM Integration}
% ===================================================================

The \texttt{SLMInterpreter} and \texttt{SLMNarrator} call a locally
running ollama instance over HTTP.

\begin{lstlisting}[title={ollama request (CommandInterpreter)}]
POST http://localhost:11434/api/generate
{
  "model": "phi3",
  "prompt": "...",
  "stream": false,
  "format": "json"
}
\end{lstlisting}

\begin{center}
\begin{tabular}{ll}
\toprule
\textbf{Property} & \textbf{Value} \\
\midrule
Default model    & phi3 (phi-3-mini, $\approx$2.3~GB) \\
Fallback         & \texttt{RulesInterpreter} if ollama unavailable \\
Context window   & 4\,096 tokens (sufficient for room context) \\
Response format  & JSON for \texttt{SLMInterpreter}; plain text for \texttt{SLMNarrator} \\
Timeout          & 5~s (\texttt{SLMInterpreter}); 10~s (\texttt{SLMNarrator}) \\
\bottomrule
\end{tabular}
\end{center}

\texttt{SLMInterpreter} system prompt provides: current room
description, available exits, visible items, and a compact action
schema.  The model returns a JSON action object that the engine
consumes directly.

\texttt{SLMNarrator} system prompt provides: room seed text from
scenario YAML plus the resolved event type and outcome.  The model
returns 1--3 sentences of atmospheric prose.

% ===================================================================
\section{Biff Extension Points}
% ===================================================================

The engine is party-ready from day one.  No redesign is required
when Biff multi-player (Phase~5) is added.

\begin{itemize}
  \item \texttt{GameState.Party} is always a \texttt{[]Character}
    slice; single-player sets \texttt{len(Party) == 1}.
  \item \tool{move}, \tool{attack}, \tool{flee}, \tool{defend} all
    accept an optional \field{character\_id}; omitting it targets
    the primary character.
  \item Combat turn order already handles multiple actors.
  \item Save files carry a \field{party} array (empty until Biff).
\end{itemize}

When Biff is added:

\begin{itemize}
  \item Each Biff participant controls one character in the party.
  \item \texttt{PostToolUse} hooks notify party members of
    state-change events.
  \item The DM narrates for the full party.
  \item Turn tokens are passed via Biff \cmd{/write} messages.
\end{itemize}

% ===================================================================
\section{Implementation Phases}
% ===================================================================

\begin{center}
\begin{tabular}{lp{10cm}}
\toprule
\textbf{Phase} & \textbf{Deliverable} \\
\midrule
\textbf{1 — Engine Core}
  & Go engine (\texttt{GameState}, \texttt{Character},
    \texttt{DungeonMap}, \texttt{Inventory}, \texttt{Combat});
    all three interfaces + all implementations except
    \texttt{LLMInterpreter}/\texttt{LLMNarrator};
    \texttt{CLIRenderer} and \texttt{LuxRenderer};
    \cmd{dungeon serve}, \cmd{dungeon solo}, \cmd{dungeon headless};
    scenario YAML schema;
    \texttt{classic-fantasy-dungeon.yaml} migration;
    save/load: \texttt{.dungeon/saves/<slot>.json}. \\
\midrule
\textbf{2 — DM Mode + Lux HUD}
  & \texttt{LLMInterpreter} + \texttt{LLMNarrator} (MCP callbacks);
    SKILL.md rewritten for DM role;
    full Lux HUD (map canvas, HP/MP bars, nav buttons, narration log);
    \cmd{dungeon dm} command;
    \cmd{/dungeon:create} --- interactive scenario generation. \\
\midrule
\textbf{3 — Combat System}
  & Turn-based combat: initiative, weapon dice, AC, conditions;
    enemy AI patterns from scenario YAML;
    Lux combat panel;
    XP, leveling, loot rolls;
    spells (Mage/Priest). \\
\midrule
\textbf{4 — Rich World}
  & Shops, traps (DEX save), cursed items, locked doors;
    multiple bundled scenarios;
    DM-generated richer encounters with backstory and puzzles. \\
\midrule
\textbf{5 — Biff Multi-player}
  & Party management (one character per Biff participant);
    turn tokens via Biff \cmd{/write};
    shared save state;
    DM narrates for the full party. \\
\bottomrule
\end{tabular}
\end{center}

% ===================================================================
\section{Open Questions}
% ===================================================================

\begin{enumerate}[label=\textbf{Q\arabic*}.]
  \item \textbf{mcp-proxy dependency.}
    DM mode (Phase~2) should use direct stdio MCP initially and
    migrate to mcp-proxy when that binary ships.  Or should Phase~2
    wait?  The daemon API is identical either way.
  \item \textbf{SLM model selection.}
    phi-3-mini ($\approx$2.3~GB), smolLM-1.7B ($\approx$1.0~GB),
    or mistral-7b-instruct ($\approx$4.1~GB)?  Evaluate against
    command-classification accuracy on a dungeon command corpus.
  \item \textbf{Go MCP SDK.}
    \pkg{github.com/mark3labs/mcp-go} versus a minimal hand-rolled
    stdio JSON-RPC layer.  The latter is $\approx$200 lines and avoids
    a dependency; the former is more complete but less battle-tested.
  \item \textbf{Dice notation parser.}
    Use a third-party Go library (e.g.\ \pkg{dicerobot/diceparser})
    or hand-roll?  The required subset (\texttt{NdM+K}) is simple
    enough to implement in $\approx$50 lines.
  \item \textbf{Permadeath default.}
    Scenario YAML can configure \texttt{death: permadeath} or
    \texttt{death: respawn}.  What should the built-in default be?
    Wizardry is permadeath; Zork is not.
\end{enumerate}

% ===================================================================
\section{Test and Verification Architecture}
\label{sec:testing}
% ===================================================================

\subsection{Philosophy}

The single most important property of the test suite is that it tells
you immediately and unambiguously whether a change broke something.
Three principles follow from this:

\begin{enumerate}[label=\textbf{P\arabic*}.]
  \item \textbf{Determinism is the invariant.}
    The engine (Section~\ref{sec:engine}) must always produce the
    same output for the same input.  Every test that touches engine
    logic asserts a deterministic outcome.  No probabilistic
    assertions, no ``roughly correct.''
  \item \textbf{Interface seams are injection points.}
    \iface{CommandInterpreter}, \iface{Narrator}, and
    \iface{Renderer} (Section~\ref{sec:interfaces}) are the only
    coupling points to external systems.  All tests above unit level
    use test doubles at these seams.  No real LLM, SLM, or Lux
    instance is ever required to run CI.
  \item \textbf{Headless mode is the CI workhorse.}
    \mode{headless} uses \texttt{RulesInterpreter} +
    \texttt{TemplateNarrator} + \texttt{CLIRenderer} — all
    zero-dependency implementations.  It runs complete game sessions
    programmatically and is the primary vehicle for integration and
    end-to-end tests.
\end{enumerate}

\subsection{The Testing Pyramid}

\begin{center}
\begin{tabular}{llrp{5.5cm}}
\toprule
\textbf{Layer} & \textbf{Trigger} & \textbf{Target} & \textbf{Coverage} \\
\midrule
Unit           & Every commit & $<$5~s   & Pure engine functions, parsers, \texttt{RulesInterpreter}, \texttt{TemplateNarrator} \\
Integration    & Every commit & $<$30~s  & Interface compositions, mock ollama, mock Lux, daemon routing \\
End-to-End     & Every PR     & $<$2~min & Real \cmd{dungeon} subprocess; MCP wire smoke; save round-trips \\
Acceptance     & Release      & $<$10~min & YAML game scripts through \cmd{dungeon headless}; full adventures \\
\bottomrule
\end{tabular}
\end{center}

All four layers require zero external services.  The full unit and
integration pass runs in under 30~seconds.  If any target is
breached, investigate before adding more tests --- the usual
culprits are sleeping tests, real socket dials, or oversized
fixture reads.  Slow tests get skipped; skipped tests catch nothing.

\subsection{Layer 1 --- Unit Tests}

\textbf{Location:} \texttt{internal/engine/*\_test.go},
\texttt{internal/interpreter/*\_test.go},
\texttt{internal/narrator/*\_test.go}

\textbf{Style:} Pure functions only.  Table-driven.  No goroutines,
no sockets, no filesystem access outside \texttt{testdata/}.

\subsubsection{Engine Core}

\begin{center}
\begin{tabular}{lp{5cm}p{4cm}}
\toprule
\textbf{Package} & \textbf{What to test} & \textbf{Key cases} \\
\midrule
\pkg{engine/character}
  & Stat bonuses, HP/MP calculation, condition apply/expire
  & All six stat values; boundary cases (8, 18); stacking conditions \\
\pkg{engine/combat}
  & Initiative order, attack roll, damage, AC reduction, flee DEX check
  & Hit/miss at boundary (roll == AC); multi-actor initiative \\
\pkg{engine/inventory}
  & Pick up/drop weight check, equip slot conflict, use-item effect
  & Full inventory rejection; slot already occupied; consumable self-removes \\
\pkg{engine/map}
  & Move through open/locked/hidden doors; fog reveal; stairway traversal
  & Locked without key; hidden pre/post search; one-way return block \\
\pkg{engine/leveling}
  & XP threshold crossing, stat delta per class
  & Level 1$\to$2, 9$\to$10; class-specific thresholds \\
\pkg{engine/combat/ai}
  & Enemy AI pattern execution
  & \texttt{aggressive} always attacks; \texttt{cautious} flees at 30\%~HP; \texttt{scripted} follows sequence \\
\pkg{engine/save}
  & JSON round-trip fidelity
  & Full \texttt{GameState} marshal/unmarshal deep-equal; version field present \\
\bottomrule
\end{tabular}
\end{center}

\textbf{Target coverage:} \texttt{internal/engine} $\geq$ 90\%
statement coverage; enforced by CI.

\subsubsection{Parsers and Formats}

\begin{center}
\begin{tabular}{lp{9cm}}
\toprule
\textbf{Package} & \textbf{What to test} \\
\midrule
\pkg{scenario/parser}
  & Valid minimal YAML loads cleanly; missing required field $\to$ typed
    error; unknown field $\to$ warning, not error; enemy template
    reference resolves; broken room exit reference $\to$ error \\
\pkg{dice/parser}
  & \texttt{1d6}, \texttt{2d6+3}, \texttt{1d20-1}; boundary rolls
    (min, max); invalid notation $\to$ error \\
\pkg{save/loader}
  & Load well-formed save; load save with unknown fields (forward
    compat); wrong \texttt{version} $\to$ typed migration error \\
\bottomrule
\end{tabular}
\end{center}

\subsubsection{RulesInterpreter}

Exhaustive verb/noun coverage.  This interpreter calls no model.

\begin{lstlisting}[title={RulesInterpreter verb table (representative)}]
"go north" / "n"       -> {action: move, direction: N}
"attack goblin" / "a"  -> {action: attack, target: "goblin"}
"pick up sword"        -> {action: pick_up, item: "sword"}
"search walls"         -> {action: search, target: "walls"}
"i" / "inventory"      -> {action: get_inventory}
"rest"                 -> {action: rest}
(unrecognised input)   -> {action: unknown, raw: "..."}  // not an error
\end{lstlisting}

\subsubsection{TemplateNarrator}

Table-driven: event type + room seed $\to$ expected substrings.

\begin{lstlisting}[title={TemplateNarrator test cases (representative)}]
{event: moved,        room_seed: "reeks of smoke"} -> contains seed text
{event: combat_start, enemies: ["goblin x3"]}      -> contains enemy count
{event: item_found,   item: "rusty_key"}           -> contains item name
{event: level_up,     new_level: 4}                -> contains "level 4"
\end{lstlisting}

\subsection{Layer 2 --- Integration Tests}

\textbf{Location:} \texttt{internal/**/integration\_test.go},
\texttt{cmd/**/integration\_test.go}

\textbf{Build tag:} \texttt{-tags integration}

\textbf{Style:} Real implementations wired together; external
services replaced with in-process fakes.  May use goroutines and
in-memory channels.  No subprocess spawning; no real network outside
\texttt{net/http/httptest}.

\subsubsection{Engine State Machine Sequences}

Multi-step action sequences exercising real engine state.  These
catch regressions that unit tests miss because they only test
individual functions.

\begin{lstlisting}[title={Representative integration sequences}]
# Locked door flow
new_game -> move(S) -> move(E) [blocked: locked] ->
pick_up(rusty_key) -> move(E) [now unlocked]
Assert: room progression correct; key consumed

# Full combat cycle
new_game -> move(S) -> {combat begins} ->
attack x N -> end_combat
Assert: HP decrements correctly; XP awarded; loot in inventory

# Level-up crossing
Accumulate XP across several combats
Assert: level increased; max_HP increased per class table

# Save/load isolation
new_game -> move(S) -> save -> mutate -> load
Assert: loaded state equals pre-mutation state
\end{lstlisting}

\subsubsection{Mock ollama HTTP Server}

\texttt{net/http/httptest.NewServer} serves canned fixture JSON.
Tests assert that:

\begin{itemize}
  \item \texttt{SLMInterpreter} correctly parses a JSON action from
    a canned response.
  \item HTTP timeout $\to$ automatic fallback to
    \texttt{RulesInterpreter} (not a crash).
  \item Non-200 response $\to$ fallback.
  \item Malformed JSON $\to$ fallback (not panic).
\end{itemize}

\subsubsection{Mock Lux MCP Server}

An in-process fake that records all \tool{show()} and
\tool{update()} calls and can inject synthetic interaction events
(button presses) via a channel.  Tests assert that
\texttt{LuxRenderer}:

\begin{itemize}
  \item Calls \tool{show()} on scene transitions (room entry, combat
    start/end); calls \tool{update()} (not \tool{show()}) for
    incremental HP, MP, and log changes.
  \item Routes \tool{recv()} button events to the engine as
    \texttt{InputEvent}s without passing through the Narrator.
  \item Does not call \tool{show()} for every HP tick (regression
    guard against the performance anti-pattern).
\end{itemize}

\subsubsection{Daemon Session Routing}

Start the daemon in-process with two fake client connections.
Assign one the DM identity; the other a player identity.

\begin{itemize}
  \item DM-privileged tools succeed from the DM client; fail from
    the player client.
  \item After any player action, \texttt{tools/list\_changed} is
    pushed to the DM client only.
  \item Client disconnection does not crash the daemon.
  \item Two concurrent game sessions are fully isolated.
\end{itemize}

\subsection{Layer 3 --- End-to-End Tests}

\textbf{Location:} \texttt{e2e/}

\textbf{Build tag:} \texttt{-tags e2e}

\textbf{Style:} Black-box subprocess tests via \texttt{os/exec}.
Requires \cmd{dungeon} binary pre-built.  Asserts observable
outputs: exit codes, JSON responses, save files on disk.

\begin{center}
\begin{tabular}{lp{4.5cm}p{3.5cm}}
\toprule
\textbf{Test} & \textbf{Command} & \textbf{Assert} \\
\midrule
Help exits clean
  & \cmd{dungeon --help}
  & Exit~0; usage text present \\
List scenarios
  & \cmd{dungeon list-scenarios}
  & JSON array; $\geq$1 entry \\
Headless new game
  & \cmd{dungeon headless --scenario minimal}
  & Exit~0; initial state JSON \\
Scripted run
  & \cmd{dungeon headless --script minimal-run.yaml}
  & All steps pass; final state matches \\
Save round-trip
  & new\_game $\to$ save $\to$ kill $\to$ load
  & Save file created; state restored \\
MCP wire smoke
  & spawn \cmd{dungeon serve}; call each tool once
  & Valid JSON for all tools; no unexpected stderr \\
Cross-mode save
  & headless save $\to$ \cmd{dungeon solo --dry-run}
  & \texttt{GameState} matches expected values \\
\bottomrule
\end{tabular}
\end{center}

\subsection{Layer 4 --- Acceptance Tests}

\textbf{Location:} \texttt{e2e/acceptance/}

\textbf{Why headless?}  Acceptance tests must be deterministic.
\mode{headless} uses \texttt{RulesInterpreter} (no model) and
\texttt{TemplateNarrator} (no model).  The engine outcome is
identical for the same scenario and input sequence on every run.

\subsubsection{Game Script Format}

Acceptance tests are expressed as YAML game scripts.  This makes
them readable and reviewable independently of Go test infrastructure.

\begin{lstlisting}[title={testdata/scripts/complete-run.yaml (excerpt)}]
scenario: minimal
character: {name: TestHero, class: fighter}
steps:
  - input: "go south"
    expect: {current_room: goblin_lair, narration_contains: "goblin"}

  - input: "attack"
    expect: {combat_active: true, enemy_hp_reduced: true}

  - input: "attack"
    repeat_until: {combat_active: false}
    max_iterations: 20
    expect: {xp_gained: true}

  - input: "save"
    expect: {save_file_exists: true}
\end{lstlisting}

\subsubsection{Bundled Acceptance Scenarios}

\begin{center}
\begin{tabular}{lp{9cm}}
\toprule
\textbf{Script} & \textbf{Coverage} \\
\midrule
\texttt{minimal-run.yaml}  & Two rooms, one combat, one item, save/load \\
\texttt{combat-full.yaml}  & All combat actions: attack, defend, use item, flee, cast spell \\
\texttt{leveling.yaml}     & XP crosses level threshold; stat increase verified \\
\texttt{locked-door.yaml}  & Find key, unlock door, search reveals hidden passage \\
\texttt{permadeath.yaml}   & Character death in permadeath scenario; game-over state asserted \\
\bottomrule
\end{tabular}
\end{center}

\subsection{Test Doubles Reference}

All fakes live in \texttt{internal/testutil/} and implement the
same Go interfaces as real implementations.

\begin{center}
\begin{tabular}{lp{5cm}l}
\toprule
\textbf{External System} & \textbf{Test Double} & \textbf{Used in} \\
\midrule
LLM (Claude)
  & \texttt{FakeLLMInterpreter} --- canned action from fixture
  & Integration \\
LLM (Claude)
  & \texttt{FakeLLMNarrator} --- fixture string
  & Integration \\
SLM (ollama)
  & \texttt{httptest.NewServer} serving fixture JSON
  & Integration \\
Lux MCP
  & \texttt{FakeLuxServer} --- records calls, injects events
  & Integration \\
Daemon socket
  & In-process fake transport
  & Integration \\
CLI stdin/stdout
  & \texttt{bytes.Buffer}
  & Integration \\
\bottomrule
\end{tabular}
\end{center}

\subsection{Verification Beyond Tests}

\subsubsection{Scenario Validation}

\cmd{dungeon validate <file.yaml>} validates a scenario against the
engine's schema and resolves all room and template references.  Run
in CI as a pre-check on any changed \texttt{scenarios/*.yaml}.
This is the safeguard for DM-authored content --- broken references
are caught before a player encounters them at runtime.

\subsubsection{MCP Schema Contract}

The MCP tool definitions are serialised to
\texttt{testdata/mcp-schema.json} by a generator.  CI diffs the
generated schema against the committed version:

\begin{lstlisting}
go run ./cmd/dump-mcp-schema > /tmp/schema.json
diff testdata/mcp-schema.json /tmp/schema.json
\end{lstlisting}

Any accidental API change --- renamed tool, added required argument,
changed return field --- fails the build before any consumer breaks.

\subsubsection{Race Detection}

\texttt{go test -race} is mandatory on CI for all packages that
touch the daemon's goroutine model.  The daemon handles concurrent
MCP connections; a data race here produces silent state corruption,
which is the hardest class of bug to diagnose in a running game
session.

\subsubsection{Save Format Forward Compatibility}

When the save format changes, a migration test loads every fixture
in \texttt{testdata/saves/} (pinned to older format versions) and
asserts clean migration to the current format.  Old saves must
never silently corrupt.

\subsection{CI Configuration}

\begin{lstlisting}[title={.github/workflows/test.yml (target layout)}]
unit-and-integration:
  - go test -race -count=1 ./...
  - go test -race -tags integration -count=1 ./...
  - go test -cover -coverprofile=coverage.out ./internal/engine/...
  - go tool cover -func=coverage.out   # fail if engine < 90%

e2e:
  - go build -o dungeon ./cmd/dungeon
  - go test -tags e2e ./e2e/...

acceptance:
  runs-on: main and release branches only
  - go build -o dungeon ./cmd/dungeon
  - go test -tags acceptance -timeout 10m ./e2e/acceptance/...
\end{lstlisting}

\subsection{What We Do Not Test Automatically}

\begin{center}
\begin{tabular}{lp{8cm}}
\toprule
\textbf{Thing} & \textbf{Reason} \\
\midrule
LLM narration quality     & Non-deterministic; evaluated by human playtest \\
SLM command accuracy      & Requires real model; evaluated by model eval harness, not CI \\
Lux visual rendering      & ImGui pixel output; evaluated by human inspection \\
mcp-proxy (pre-ship)      & Does not exist yet; contract test stubs only \\
Balance and pacing        & Game design; evaluated by playtesting \\
\bottomrule
\end{tabular}
\end{center}

\end{document}
